% \documentclass[openacc]{rsproca_new}
\documentclass[a4paper]{article}
% \titlehead{Research}

\usepackage[margin=1in]{geometry}
\usepackage{hyperref}
\usepackage{color}
\usepackage{graphicx}
\usepackage{amsmath}
\usepackage{amssymb}
\numberwithin{equation}{section}
\usepackage[backend=bibtex, style=numeric-comp, sorting=none]{biblatex}
\DeclareFieldFormat{labelnumberwidth}{[#1] }
\addbibresource{main.bib}

\begin{document}

\title{Statistical mechanics of Hamiltonian curl-force systems}

\author{%%%% Author details
Omer Chor$^{1}$, Alja\v{z} Godec$^{2}$ and Oren Raz$^{1}$}

%%%%%%%%% Insert author address here
\date{
$^{1}$Department of Physics of Complex Systems, Weizmann Institute of Science, Rehovot, 76100, Israel
\\
$^{2}$ \textcolor{red}{Add Alja\v{z}'s affiliation}\\[2ex]
\today}

\maketitle

%%%% Subject entries to be placed here %%%%
% \subject{xxxxx, xxxxx, xxxx}

%%%% Keyword entries to be placed here %%%%
% \keywords{xxxx, xxxx, xxxx}

%%%% Insert corresponding author and its email address}
% \corres{Insert corresponding author name\\
% \email{xxx@xxxx.xx.xx}}

% \date{\today}

\begin{abstract}
We investigate the statistical mechanics of a particle subject to curl-forces---reversible, velocity independent non-conservative forces with non-vanishing curl---that can be described by a Hamiltonian with an anisotropic kinetic energy term. We show that generally they do not relax to the corresponding Boltzmann distribution when coupled to a heat bath, and therefore the Hamiltonian structure cannot be used to easily find their nonequilibrium steady state. By applying a canonical transformation, the system can be mapped to a system without curl-forces that is instead coupled to multiple heat baths at different temperatures. For quadratic potentials, this mapping recovers the underdamped Brownian gyrator, providing a new perspective on the origin of its rotational currents as the result of deterministic, non-conservative forces. The same analysis is also applied to another variation of the Brownian gyrator, demonstrating its generality.
\end{abstract}

\maketitle

\section{Introduction}
\label{sec:Intro}
Equilibrium statistical mechanics is by now a mature field. One of its most remarkable results is the universal structure of the equilibrium distribution: any system that can exchange heat with a (single) large enough thermal bath eventually relaxes into an equilibrium distribution, the Boltzmann distribution, that is completely determined by the system's Hamiltonian and the bath's temperature. At the microscopic level, the stability of this distribution is guaranteed by the principle of detailed balance, which ensures the absence of net probability currents \cite{seifert2012stochastic}.
% where the probability to find the system in a microstate $\mu_i$ is given by the Boltzmann distribution, $P(\mu_i) = e^{-H(\mu_i)/k_BT}/Z$, where $H(\mu_i)$ is the value of the Hamiltonian---the energy functional of the system---at this microstate, $T$ is the bath temperature, $k_B$ is the Boltzmann constant and $Z$ is the partition function, which is the corresponding normalization.
% Unfortunately, we do not know of an analog result for systems that are not in thermal equilibrium. In some nonequilibrium scenarios, e.g. in periodically driven system or glassy dynamics, the system does not relax into a steady state, and the probability to find the system in a given microstate changes over time or strongly depends on the exact initial condition. But even in the simplest type of nonequilibrium, where detailed balance is violated but yet the system relaxes into a NESS (nonequilibrium steady state), we have no simple formula for this steady state, and finding it is a difficult problem that to the best of our knowledge has to be solved case by case.

Unfortunately, no such universal principle exists for systems that are not in thermal equilibrium. Some nonequilibrium scenarios, such as glassy or periodically driven systems, may never reach a stationary state. But even the seemingly simpler case of a Non-Equilibrium Steady State (NESS) poses a significant theoretical challenge \cite{sasa2006steady}. In a NESS, the violation of detailed balance is manifested in persistent probability fluxes \cite{seifert2012stochastic, sasa2006steady, zia2007probability}, such that the stationary distribution alone does not uniquely describe the system \cite{zia2007probability}. As a consequence, there is no direct analog of the Boltzmann distribution for a NESS, and exact solutions typically require a case-by-case analytical approach.

There are many physical scenarios where the detailed balance condition does not hold and the system is driven out of equilibrium: some degrees of freedom (DoF) might be driven microscopically as in active matter \cite{marchetti2013hydrodynamics, bechinger2016active}; feedback control protocols \cite{seifert2012stochastic, sagawa2012nonequilibrium, parrondo2015thermodynamics, goerlich2025experimental} and time dependent driving \cite{jung1993periodically, schmiedl2008efficiency, raz2016mimicking,busiello2018similarities, koyuk2020thermodynamic}; a system connected to two or more baths with different temperatures, chemical potentials or voltages \cite{zia2007probability, derrida2007non}; or the system can be influenced by some external non-conservative force \cite{seifert2012stochastic, sasa2006steady, chernyak2006path}. The two latter cases are the main focus of this manuscript. 

The study of systems coupled to multiple reservoirs originates in the very foundation of classical thermodynamics, where the heat engine operating between a hot and a cold reservoir serves as the paradigmatic setup studied by Carnot \cite{carnot1872reflexions}. Clausius's formulation of the second law of thermodynamics also concerns heat flows between reservoirs \cite{clausius1879mechanical}. Later on, the transport associated with gradients in the value of thermodynamic quantities was considered by Onsager in his seminal works on reciprocal relations \cite{onsager1931reciprocal}. Systems coupled to multiple baths remain of great interest to date, where typically a microscopic model is considered instead of the earlier macroscopic descriptions. These include boundary-driven models such as asymmetric simple exclusion processes (ASEP) \cite{krug1991boundary, derrida1998exactly, derrida2007non} or oscillator chains \cite{kipnis1982heat, eckmann1999non}. More recently, there has also been great interest systems where different DoF are coupled to a different heat bath within the framework of stochastic thermodynamics \cite{seifert2012stochastic}. Notable examples are Brownian motors and refrigerators \cite{van2004microscopic, van2006brownian, filliger2007brownian}, which are a modern take on the classical macroscopic heat engine. These ideas have been generalized and studied extensively both theoretically and experimentally \cite{ dotsenko2013two, mancois2018two, bae2021inertial, baldassarri2020engineered,cerasoli2018asymmetry, cerasoli2022spectral, argun2017experimental, martins2026role, chang2021autonomous, dos2021stationary, pietzonka2018universal, lin2022stochastic, movilla2021energy, abdoli2022tunable, chiang2017electrical, ciliberto2013heat, PhysRevE.92.032118, abdoli2025quadrupolar, visco2006work, berut2014energy, berut2016stationary, fogedby2011bound, fogedby2017minimal, leighton2024information, netz2020approach} in the past two decades. In these systems, the simultaneous coupling to several reservoirs imposes incompatible equilibrium conditions; consequently, the system cannot relax to a standard Boltzmann distribution, but settles instead into a NESS characterized by persistent currents and non-vanishing entropy production. Given the well-established nature of these systems, it is particularly compelling to determine when seemingly distinct driving mechanisms, such as the non-conservative forces discussed below, can be formally mapped onto this reservoir-driven paradigm for nonequilibrium thermodynamics.

In principle, all the fundamental forces in nature are conservative, and can be derived from a potential energy function. However, in cases where a complete microscopic description is unavailable for some parts of the total system, it  is often necessary to approximate the dynamic through non-conservative forces. For example, drag forces are often introduced to account for the influence of a medium on the motion of a particle. Generally, non-conservative forces can be \emph{dissipative}, which means that the dynamics of the system are not reversible and do not conserve phase space volume. Such non-conservative forces do not necessarily violate the detailed balance condition (and evidently, they are necessary for the equilibration of an excited system). However, there also exist reversible non-conservative forces, which conserve phase space volume but are not derived from a potential. Specifically, we are interested in the case of force fields $\mathbf{F}(\mathbf{x})$ that depend only on the position $\mathbf{x}$ of the particles but not on their velocities, and are not the gradient of any scalar potential. In 3 dimensions, this implies that they have a non-vanishing curl, $\nabla \times \mathbf{F}(\mathbf{x})\neq0$. Following Berry and Shukla, these are referred to as \emph{curl-forces} \cite{berry2012classical}. 

Such curl-forces are usually associated with external forcing, which typically drives a thermal system out of equilibrium. These forces are not merely theoretical constructs: they can be generated in the lab, for example through an optical force field \cite{curtis2003structure, roichman2008influence,roichman2008optical, sun2009brownian, wu2009direct, liberal2013near, sukhov2017non, volpe2023roadmap}, or using feedback traps \cite{jun2012virtual}. 
A series of recent papers studied various properties of classical and quantum systems with curl-forces 
\cite{berry2012classical, berry2015hamiltonian, berry2016curl, berry2023quantum, berry2024quantising, berry2025six}. 
In particular, it was shown that certain two-dimensional systems with curl-forces can be described using Hamiltonian formalism. In this description, however, the Hamiltonian function differs from the canonical structure of kinetic plus potential energies \cite{berry2015hamiltonian}. The fact that some curl-forces can be described by a (non-standard) Hamiltonian raises a natural question: can the NESS of such systems be described by the Boltzmann distribution corresponding to their Hamiltonian? 

In this work, we show that unfortunately this is not the case: the Hamiltonian corresponding to a curl-force field cannot be used to easily find the steady state of systems driven out of equilibrium. However, we demonstrate how curl-forces can be used to give a clear qualitative description for currents arising in other, more obscure settings.
We begin in Sec.~\ref{sec:HCF} by a short introduction on Hamiltonian curl-forces. In Sec.~\ref{sec:StochDyn} we discuss the statistical mechanics of such systems when coupled to a heat bath, and explain why they do not relax to a Boltzmann distribution despite their Hamiltonian structure. We then consider the Klein-Kramers equation describing the stochastic evolution of the system.
In Sec.~\ref{sec:Transformation}, we present our main result: using a specific canonical transformation, we show that there is a mapping between thermal systems with Hamiltonian curl-forces and systems without any curl-forces that are instead simultaneously coupled to several heat baths at different temperatures, which is a different class of nonequilibrium systems. We study two examples of such systems in Sec.~\ref{Sec:Applications}, both featuring rotational currents solely due to the coupling to two separate heat baths, without any curl-forces. By analyzing $\nabla \times \mathbf{F}$ in the transformed, curl-force description of the same systems, we are able to provide a simple mechanical explanation for the origin of these currents. Finally, in Sec.~\ref{sec:discussion} we summarize our results and provide an outlook with open questions and possible generalizations.


\section{Hamiltonian curl-forces}

\label{sec:HCF}
In this section we give a short reminder on Hamiltonian curl-forces, following Berry and Shukla \cite{berry2015hamiltonian}. For simplicity, we first discuss systems with two DoF, and then extend to systems with more DoF.

The standard Hamiltonian of a classical system (without a magnetic field) is of the form 
\begin{equation}\label{eq:Classical_Hamiltonian}
    H(\mathbf p,\mathbf x) = \frac{\mathbf{p}^2}{2m} + U(\mathbf x)
\end{equation}
where $\mathbf p=(p_{1},p_{2})$ is the canonical momentum conjugate to the cartesian coordinates $\mathbf x = (x_1,x_2)$\footnote{In principle, one can use any set of canonical coordinates and momenta, not necessarily cartesian coordinates. However, in the non cartesian case it is less trivial to identify whether the force field is velocity independent. Therefore, we assume here cartesian coordinates for simplicity.}, and $m$ is the particle's mass. The equations of motion of the system are
\begin{eqnarray} \label{eq:EqOfMotion}
    \dot x_i = \frac{\partial H}{\partial p_i} = \frac{p_i}{m} \qquad 
    \dot p_i=-\frac{\partial H}{\partial x_i} = -\frac{\partial U(\mathbf x)}{\partial x_i} 
\end{eqnarray}
The force is defined as 
\begin{equation} \label{eq:NewtonForce}
    \mathbf F \equiv m\ddot{\mathbf x} = -\nabla U(\mathbf x),
\end{equation}
which is necessarily conservative, as it is a gradient of the potential. 

One can, in principle, consider Hamiltonians that are not of the form Eq.~\ref{eq:Classical_Hamiltonian}. The force in this case is still defined as $\mathbf F = m \ddot{\mathbf x}$, but the force field $\mathbf F$ might be velocity dependent, and not only a function of position. Berry and Shukla showed \cite{berry2015hamiltonian} that  under mild assumptions, the only Hamiltonian structure that leads to a velocity independent force field, is 
\begin{eqnarray} \label{eq:CurlHamiltonians}
    H(\mathbf p,\mathbf x) = \frac{1}{2m}\mathbf p^T \mathcal{A} \mathbf p + U(\mathbf x),
\end{eqnarray}
where $\mathcal{A}$ is a symmetric, position independent matrix. In this case,
\begin{equation}
\mathbf F \equiv m\ddot{\mathbf{x}} = -\mathcal{A}\nabla U(\mathbf x).
\label{eq:HCF-def}
\end{equation}
Without loss of generality, let us work in the (cartesian) basis where $\mathcal{A}$ is diagonal, 
\begin{eqnarray}
    \mathcal{A} = \begin{pmatrix}
         \alpha & 0 \\
         0 & \gamma
    \end{pmatrix} ,
\end{eqnarray}
where $\alpha$ and $\gamma$ are unitless constants.
The curl of the force field is given by
\begin{equation}\label{eq:berry_curl}
    \nabla\times\mathbf F = \left( \alpha - \gamma \right) \partial_{x_1}\partial_{x_2}U 
\end{equation}
and therefore the force field is not curl free, unless either $\alpha=\gamma$ or $U(\mathbf x)$ is everywhere additively separable in the coordinates $\mathbf x$, namely $U(\mathbf x) = f(x_1)+g(x_2)$ for some functions $f$ and $g$. 

The above argument for two DoF can be generalized to a system with more DoF as follows: instead of $\nabla \times \mathbf F$, consider the matrix $(\nabla \mathbf F)_{ij}\equiv \partial_iF_j$. From Eq.(\ref{eq:HCF-def}) it follows that $\nabla \mathbf F = -\mathcal{A}\nabla\nabla U(\mathbf x)$ where $(\nabla\nabla U)_{i,j} = \partial_{x_i}\partial_{x_j}U$, namely the Hessian of $U$. If it is symmetric, then by Poincare lemma $\mathbf F = \nabla \phi (\mathbf x)$ for some scalar potential $\phi(\mathbf x)$ (we assume a simply connected manifold). Therefore, $\mathbf F$ is not a gradient of a potential only if $\mathcal{A}\nabla\nabla U\neq \nabla\nabla U \mathcal{A}$ (note that both the Hessian and $\mathcal{A}$ are symmetric matrices). This means that the force field $\mathbf F$ is not conservative --- it is not a gradient of a potential --- when the matrix $\mathcal{A}$ and the Hessian of $U$ do not commute, namely $[\mathcal{A},\nabla\nabla U(\mathbf{x})]\neq 0$ at some $\mathbf{x}$.  

\section{Statistical mechanics}
\label{sec:StochDyn}
Equilibrium systems that can exchange heat with a big enough thermal bath eventually relax to a Boltzmann distribution, where the probability of the system to be in a microstate $\mu_i$ is given by $P(\mu_i)=e^{-H(\mu_i)/k_B T}/Z$. Here $H(\mu_i)$ is the energy of the system at microstate $\mu_i$, $T$ is the bath temperature, $k_B$ is Boltzmann's constant, and $Z$ is the partition function, which is the corresponding normalization. This famous result due Boltzmann is based on conservation of the total energy, 
\begin{equation}\label{eq:total_energy}
    H_{tot} = H_{sys}(\mathbf p_{sys},\mathbf x_{sys})+H_{bath}(\mathbf p_{bath},\mathbf x_{bath}) +  U_{int}(\mathbf x_{sys},\mathbf x_{bath}),
\end{equation}
where $H_{sys}(\mathbf p_{sys},\mathbf x_{sys})$ is the Hamiltonian of the system with DoF $(\mathbf p_{sys},\mathbf x_{sys})$, $H_{bath}(\mathbf p_{bath},\mathbf x_{bath})$ is the energy of the thermal bath and is only a function of the bath DoF, and $U_{int}(\mathbf x_{sys},\mathbf x_{bath})$ is the energy associated with the interaction between the the system and the bath, which depends on both sets of coordinates but not on their momenta. To get the Boltzmann distribution, it is assumed that $U_{int}(\mathbf x_{sys},\mathbf x_{bath})$ is negligible in comparison to $H_{sys}$ and $H_{bath}$. 

We are interested in systems driven by a non-conservative curl-force, whose Hamiltonian $H_{sys}$ has an anisotropic kinetic energy as in Eq.~\ref{eq:CurlHamiltonians}. If this Hamiltonian has no explicit time dependence, then for an isolated system its value is conserved, and can be viewed as `energy', in the sense that this is the conserved quantity associated with the time translation symmetry of the system. 

In order to describe the statistical mechanics of the driven system when coupled to a thermal bath, we need to make some assumptions regarding the coupling of the system to the bath. Our goal is to use the curl-force Hamiltonian formalism to describe a standard thermal system subject to non-conservative curl-forces. Essentially, if the driving is turned off, we require that the system's dynamics reduce to those of a system in thermal equilibrium.
Therefore, we assume that the bath itself and the interaction of the system with the bath are the standard equilibrium ones, such that the only source of nonequilibrium is contained within its non-standard Hamiltonian. In terms of $H_{sys}$ and the interaction energy $U_{int}$, this means that the equations of motion of the system DoF are
\begin{equation}\label{eq:newton_system_bath}
      m\ddot{\mathbf{ x}}_{sys} = -\mathcal{A}\cdot\nabla U_{sys}(\mathbf x_{sys}) -\nabla U_{int}(\mathbf x_{sys},\mathbf x_{bath}).
\end{equation}
In this case, $H_{tot} = H_{sys}+H_{bath} +  U_{int}$ is \emph{not} a conserved quantity, since it is not true anymore that the equations of motion for the system DoF follow the Hamiltonian dynamics associated with 
$H_{sys}$. Therefore, the Boltzmann distribution is no longer expected to be the steady state distribution describing the system, and in general the system will not relax into an equilibrium state. 

We note that if the system's dynamics were instead described by
\begin{equation}\label{eq:newton_system_modified}
    m\ddot{\mathbf{ x}}_{sys} = -\mathcal{A}\cdot\nabla\Big(U_{sys}(\mathbf x_{sys}) + U_{int}(\mathbf x_{sys},\mathbf x_{bath})\Big)
\end{equation}
rather than Eq.~\ref{eq:newton_system_bath}, that is, if the interaction with the bath would also scale according to the anisotropy matrix $\mathcal{A}$, then the dynamics would follow the canonical equations of motion, and as a result the total Hamiltonian would be a conserved quantity. The corresponding system would then indeed relax to the Boltzmann distribution associated with $H_{sys}$. However, that case no longer describes the dynamics of a regular particle under the influence of an external curl-force interacting with a bath, and hence it cannot be used to infer the steady state distribution of a particle subject to a non-conservative force. 

\subsection{Stochastic dynamics}
The equations of motion described in Eq.~\ref{eq:newton_system_bath} have two components: one is the potential (curl) force, $\mathbf F_{pot}=-\mathcal{A}\cdot\nabla U_{sys}$, and the other describes the interaction with the bath, $\mathbf F_{int}=-\nabla U_{int}$. In order to get a stochastic description of the particle's motion, the interaction force is usually coarse grained and described as the combination of a dissipative force and a random fluctuating force. Under the standard assumptions that the bath is large and the interaction energy is small, the assumption that the bath remains in equilibrium is still valid for a system with curl-forces, and we therefore assume that the dissipation and random force satisfy the fluctuation-dissipation relation as usual. 
We emphasize that the canonical momentum of the curl-force system is not a physically measurable quantity, and the interaction with the bath should depend on the immersed particle's velocity and not on its canonical momentum. This means that the dynamical friction must be proportional to the particle's velocity vector $\dot{\mathbf{x}}$, which is not proportional to the canonical momentum vector $\mathbf{p}$ if $\mathcal{A}$ is not the identity matrix. This ensures that when the potential is turned off, the particle diffuses like a regular particle, and is in line with the literature on Brownian motion in other cases where the canonical momentum $\mathbf{p}$ does not coincide with the inertial momentum $m\dot{\mathbf{x}}$, for example for a charged particle in a magnetic field \cite{taylor1961diffusion, kurcsunoglu1963brownian, czopnik2001brownian, simoes2005kramers, lisy2013brownian}.

With these assumptions, we can write the Klein-Kramers equation describing the underdamped evolution of the probability density function $f(\mathbf{x},\dot{\mathbf{x}})$ of the system as
\begin{equation}
    \label{eq:klein_kramers}
    \frac{\partial f(\mathbf{x}, \dot{\mathbf{x}})}{\partial t} =  \frac{\mathcal (\mathcal A\nabla_{\mathbf{x}} U)^T \cdot \nabla_{\dot{\mathbf{x}}}f}{m}-\nabla_{\mathbf{x}} f \cdot \dot{\mathbf{x}} + \frac{\xi}{m} \nabla_{\dot {\mathbf x}} \cdot(\dot {\mathbf x} f + \frac{k_B T}{m} \nabla_{\dot{\mathbf{x}}}f),
\end{equation}
where $\xi$ is the friction coefficient. Here, $\nabla_{\mathbf{x}},\nabla_{\dot {\mathbf x}}$ are the del operators with respect to the coordinates and velocities, respectively. The nonequilibrium nature of this system appears through the anisotropy matrix $\mathcal{A}$ in the convection term, and it can be verified by a direct substitution that the Boltzmann distribution corresponding the curl-force Hamiltonian $H$ is indeed not a steady-state solution of this equation. We further mention that the first two terms are compactly written as
$\{H,f\}\equiv\nabla_{\mathbf x} H \cdot \nabla_{\mathbf p}f-\nabla_{\mathbf x} f \cdot \nabla_{\mathbf p}H$,
where $\{H,f\}$ is the Poisson bracket of the curl-force Hamiltonian $H$ (Eq.~\ref{eq:CurlHamiltonians}) and the probability density distribution function $f$.

\subsection{{Steady state}}
A natural question is then the exact conditions on $U, \mathcal{A},\xi$ and $T$ such that the system has a unique, normalizable steady state. %From Eq.~\ref{eq:klein_kramers} it is clear that this question is equivalent to the exact conditions for the force field $\mathbf{F}=-\mathcal{A}\nabla_{\mathbf{x}}U$ to have a unique normalizable steady state in the Klein-Kramers equation \textcolor{red}{[Omer: it's not clear to me how this sentence is different than the previous one: clearly, we claim that the KK eq. describes the system, so if it has a steady state it describes the physical system]}. 
The exact conditions for a unique normalizable steady state in a general force field are, as far as we know, not known, though for a given $\mathbf{F}(\mathbf{x})$ this can commonly be done using the Foster--Lyapunov criterion \cite{dolbeault2018varphi}. 

From the next section on, we focus on $\mathcal{A}$ which is positive definite, such that $\alpha,\gamma>0$. Under this assumption it is clear that for bounding potentials the system has a unique normalizable steady state. However, we point out that it is possible that systems with $\alpha<0$ or $\gamma<0$ have a normalizable steady state. A trivial example is the saddle-like potential
\begin{equation}
    U(x_1,x_2)=\frac{k}{2}(x_1^2-x_2^2)
\end{equation}
combined with the anisotropy matrix
\begin{eqnarray}
    \mathcal{A} = \begin{pmatrix}
         1 & 0 \\
         0 & -1
    \end{pmatrix}.
\end{eqnarray}
In this case 
\begin{eqnarray}
    \mathbf{F} = -\mathcal{A}\cdot\nabla_{\mathbf{x}} U= -k\begin{pmatrix}
        x_1\\
        x_2
    \end{pmatrix}
\end{eqnarray}
and this force field, which corresponds to a harmonic potential with a standard kinetic energy, clearly has a normalizable steady state. 

% Note to self: definitely look into \cite{kwon2005structure, qian2001mathematical}.


\section{Canonical transformation: multiple temperatures interpretation}\label{sec:Transformation}
Assuming that $\mathcal{A}$ is positive definite, it is cleat that the kinetic energy of the system is always non-negative. This assumption, however, limits the types of curl-forces that can be described: for example, the purely solenoidal force field $\mathbf{F}=y\hat{x}-x\hat{y}$ requires that either $\alpha$ or $\gamma$ are negative.
With this assumption, a different perspective to Eq.~\ref{eq:klein_kramers} can be given by the following canonical transformation. We  define a generating function of the second kind \cite{LANDAU1976131},
\begin{equation}
    G_2(\mathbf P, \mathbf x) = \mathbf P^T \cdot \mathcal A^{-{1}/{2}}\cdot\mathbf x,
\end{equation}
where $\mathbf P$ and $\mathbf x$ are the new momenta and old coordinates, respectively. From it, we obtain the new coordinates and momenta in terms of the old ones as
\begin{equation}\label{eq:canonical_transformation}
\begin{split}
    \mathbf X &= \mathcal A^{-1/2}\mathbf x,\\
    \mathbf P &= \mathcal A^{1/2} \mathbf p,
\end{split}
\end{equation}
and the transformed Hamiltonian is
\begin{equation}
    \tilde{H}(\mathbf X, \mathbf P) = \frac{1}{2m} \mathbf P^2 + \tilde{U}\left(\mathbf X\right),
\end{equation}
where $\tilde{U}(\mathbf X) = U( \mathcal{A}^{1/2} \mathbf X) = U(\mathbf{x})$ is the potential energy in terms of the new coordinates. $\tilde{H}$ has an isotropic kinetic energy term with a standard potential term, so in principle it describes a system without any curl-forces. However, the non-conservative nature of the forces is now encoded in the way the system is coupled to the environment, which leads again to nonequilibrium dynamics. 
Applying this coordinate transformation to the Klein-Kramers equation (Eq.~\ref{eq:klein_kramers}), we obtain
\begin{align}\label{eq:CK_transformed}
    \frac{\partial f(\mathbf{X}, \dot{\mathbf{X}})}{\partial t} &= \frac{1}{m}\nabla_{\mathbf{X}} \tilde{U} \cdot \nabla_{\dot{\mathbf{X}}}f-\nabla_{\mathbf{X}} f \cdot \dot{\mathbf{X}} + \frac{\xi}{m} \nabla_{\dot {\mathbf X}} \cdot(\dot {\mathbf X} f + \frac{k_B T}{m} \mathcal{A}^{-1}\cdot\nabla_{\dot{\mathbf{X}}}f) \\
    &=\{\tilde{H}, f\} + \frac{\xi}{m} \nabla_{\dot {\mathbf X}} \cdot(\dot {\mathbf X} f + \frac{k_B T}{m} \mathcal{A}^{-1}\cdot\nabla_{\dot{\mathbf{X}}}f). \nonumber
 \end{align}
This equation can be interpreted as an underdamped motion of a particle under the influence of a potential $\tilde{U}$, where the noise magnitude at each axis is different. It is constructive to interpret 
\begin{equation}
    T_i \equiv \mathcal{A}_{ii}^{-1} T
\end{equation} 
as the temperatures of a heat reservoir coupled to the $i$-th coordinate. $T_i$ is a physically valid temperature, as it is positive under our earlier assumption that $\mathcal{A}$ is positive definite. This coupling implies that the system is generally out of equilibrium, unless either $T_i$ are all equal, or the different degrees of freedom $X_i$ are decoupled, $\partial_{X_i X_j} U = 0$, and the potential is additively separable in the coordinates. In the latter case, the force field in the single-reservoir description can be expressed as minus the gradient of a potential $\mathbf{F}=-\nabla(\alpha U_1(x) +\gamma U_2(y))$, the system's energy can be written in terms of a standard Hamiltonian $H=\frac{\mathbf{p}^2}{2m}+\alpha U_1(x) +\gamma U_2(y)$ with isotropic kinetic energy, and the system eventually relaxes to the corresponding Boltzmann distribution. This is an alternative explanation for the fact that if in the multiple-reservoirs description the DoF are decoupled, heat cannot flow between the reservoirs and each DoF equilibrates with its own bath \cite{lin2022stochastic, leighton2024information, netz2020approach}. %\textcolor{red}{[Please see what you think about these sources. While they do not explicitly state this, I think it is hinted. In the first paper, they say a non-conservative coupling force is required for heat flow between the two reservoirs. In Sivak's paper they talk about the mutual information, which indicates some coupling between the subsystems. The last paper also has a statement about the EPR (or some other noneq property) vanishing if the coupling vanishes.]}. %In principle, this statement can be generalized to more complex forms of $\mathcal{A}$, namely cases where it has off-diagonal terms, where it might be harder to tell whether the system equilibrates or not. The curl-force description provides a clean interpretation for such systems, where the condition for the system to equilibrate becomes clearer.

In summary, this canonical transformation provides a mapping between Hamiltonian curl-forces and the motion of a particle under the influence of different thermal bath connected to the various DoF, and a correspondingly stretched potential. Indeed, in Eq.~\ref{eq:CK_transformed} there is no longer any curl-force, and the nonequilibrium nature of the system arises solely from the fact that now it is coupled to multiple heat baths. %We also note that the inverse transformation can be used to analyze a non-equilibrium system that is only driven away from equilibrium by coupling different DoF to different thermal baths. 

\subsection{Transforming temperature}
The mapping defined by $G_2$ uniquely defines not only the new coordinates and momenta, but also the temperatures $T_i$ of the multiple heat baths in the transformed description of the thermal system. The canonical transformation in Eq.~\ref{eq:canonical_transformation} is invertible, so one can consider the inverse mapping of $G_2$,
\begin{equation}\label{eq:inverse_canonical_transformation}
\begin{split}
    \mathbf x &= \mathcal A^{1/2}\mathbf X,\\
    \mathbf p &= \mathcal A^{-1/2} \mathbf P.
\end{split}
\end{equation}
In a non-thermal dynamical system, this implies that the transformed system describes the same dynamics as the original system, and either descriptions can be used to analyze its properties.
It is therefore sensible to consider the inverse transformation from a thermal system with multiple heat baths and no curl-forces to a system with a single heat bath and curl-forces. This is arguably the more `useful' direction of the transformation, in the sense that it can describe the origin of a nonequilibrium steady state of thermal systems with multiple baths in terms of a deterministic, non-conservative force, which is conceptually simpler. However, applying the inverse transformation of $G_2$ to a thermal system gives rise to an ambiguity in the temperature $T$ of the single bath, as $T_i=\mathcal{A}_{ii}^{-1}T$ only defines the finite temperature $T$ up to a constant that can be absorbed in $\mathcal{A}$. Clearly, $T$ must be fixed in order to provide a consistent description of the transformed system.

We couldn't find any physical constraint for $T$ other than ensuring that it is consistent with the case where the all the multiple baths have the same temperature, which effectively reduces to a system with a single bath. This still leaves some freedom in the choice of the temperature. Therefore, the inverse transformation of $G_2$ defines a set of dual single-bath systems.

A useful choice is to set the temperature as the geometric mean of the multiple baths, $T=\left(\prod_{i=1}^d T_i\right)^{1/d}$, which implies $\det \mathcal A=1$. Since the temperatures of the multiple baths are inversely related to the components of $\mathcal{A}$, their product ($\det \mathcal{A}$) appears when computing the common denominator of quantities involving addition or subtraction of the temperatures. Then, this substitution immensely simplifies computations.

% We do this by noting that in two dimensions, $\nabla \times \mathbf{F}(\mathbf{x})$ of the single-bath description (Eq.~\ref{eq:berry_curl}) can be expressed in terms of the potential $\tilde{U}(\mathbf{X})$ of the multi-bath system as 
% \begin{equation}
%     \nabla \times \mathbf{F}(\mathbf{x}) = (\alpha-\gamma) \partial_{x_1}\partial_{x_2} U = (\alpha-\gamma)\frac{\partial_{X_1}\partial_{X_2} \tilde{U}}{\sqrt{\alpha \gamma}} .
% \end{equation}
% \st{Since generally $U$ and $\tilde{U}$ are independent of $\alpha$ and $\gamma$, we must fix them to eliminate the dependence on them in the denominator (which does not appear in the l.h.s. of the equality)}. Further noting that $\alpha=\gamma=1$ is a valid physical choice which represents the identity transformation in the case of equilibrium dynamics (when there are no curl-forces, or the temperature in each direction has the same amplitude), the above normalization rule must allow it. Therefore, we fix $\alpha \gamma=1$. Generalizing to higher dimensions, we set
% \begin{equation}\label{eq:new_temperature}
%     \det \mathcal{A} = 1, \quad \text{equivalently, }\quad T=\left(\prod_{i=1}^d T_i\right)^{1/d},
% \end{equation}
% i.e., the temperature of the single bath is the geometric mean of the temperatures of the multiple baths.

\section{Applications}
\label{Sec:Applications}
In this section we provide two examples for systems that are driven out of equilibrium by different temperatures coupled to different coordinates, and analyze their curl-force representation.

\subsection{Quadratic potential: the Brownian gyrator} 
The Brownian gyrator is one of the simplest models for a stochastic heat engine. Originally proposed by Filliger and Reimann \cite{filliger2007brownian} and further studied extensively \cite{dotsenko2013two, mancois2018two, bae2021inertial, baldassarri2020engineered, cerasoli2018asymmetry, cerasoli2022spectral, argun2017experimental, martins2026role, chang2021autonomous, dos2021stationary, movilla2021energy, pietzonka2018universal, lin2022stochastic, abdoli2022tunable, chiang2017electrical}, it describes a Brownian particle diffusing in the presence of an anisotropic harmonic potential,
\begin{equation}\tilde{U}(X, Y) = \frac{1}{2}\left( kX^2 + 2uXY + kY^2 \right),\end{equation}
with $|u|<k$, and two different random noises amplitudes in the $X$ and $Y$ directions (corresponding to temperatures $T_X$ and $T_Y$, respectively).  
Due to the harmonicity of the potential, the corresponding stochastic process is an Ornstein–Uhlenbeck (OU) process, which is exactly solvable \cite{risken1996fokker}. The non-equilibrium steady-state of this system exhibits rotational motion, whose magnitude is determined by both the anisotropy of the potential, determined by $u/k$, and the temperature difference between the two heat baths, $T_Y-T_X$ \cite{mancois2018two, bae2021inertial}. By applying the inverse transformation of $G_2$ from the two temperature case of the inertial Brownian gyrator, we offer an alternative interpretation for the source of rotational motion: it can be viewed as a result of a \emph{deterministic} curl-force acting on a particle coupled to a single heat bath. The corresponding potential in the curl-force representation is
\begin{equation}
    U(x, y) = \frac{1}{2}\left( \frac{k}{\alpha}x^2 + 2\frac{u xy}{\sqrt{\alpha \gamma}} + \frac{k}{\gamma}y^2 \right) = \frac{1}{2T}\left( {kT_X}x^2+2{u\sqrt{T_x T_y}xy}+kT_yy^2 \right),
\end{equation}
and the total force is \begin{align}
\mathbf{F}(x, y) &\equiv -\mathcal{A}\cdot\nabla_\mathbf{x} U =-(kx+\sqrt{\frac{\alpha}{\gamma}} uy)\hat{x}-(ky+\sqrt{\frac{\gamma}{\alpha}} ux)\hat{y} \\
&=-(kx+\sqrt{\frac{T_Y}{T_X}} uy)\hat{x}-(ky+\sqrt{\frac{T_X}{T_Y}} ux)\hat{y}. \nonumber
\end{align} The curl of the force field in this case is constant, given by 
$\nabla \times \mathbf{F} = u\left({\frac{\alpha-\gamma}{\sqrt{\alpha\gamma}}}\right)$, or $\nabla\times\mathbf{F}=u\left(\frac{T_Y - T_X}{\sqrt{T_X T_Y}}\right)$ in terms of the temperatures of the multiple baths $T_{X,Y}$.
It is non-vanishing unless either  $\alpha = \gamma$ and hence $T_X = T_Y$ and there is a single termal bath, or $u=0$, namely the principal axes of the potential coincides with the two different thermal baths, such that the $X$ and $Y$ coordinates become decoupled. For $T_X>T_Y$ ($\alpha<\gamma$), the curl is negative, which corresponds to the force field inducing a clockwise circulation. This result is consistent with the results obtained for the inertial Brownian gyrator \cite{mancois2018two, bae2021inertial}. Notably, this qualitative result is purely mechanical, and does not require a solution of the Fokker-Planck equation associated with the stochastic dynamics of the system. It gives an easy way to obtain and intuitive explanation for the origin and the direction of the rotation of the Brownian particle in the two-temperatures setup. 

It is instructive to compare the energetics of the two descriptions. In the Brownian gyrator, heat $\langle \dot{Q}_X \rangle = -\langle \dot{Q}_Y \rangle$ flows from the hot to the cold reservoir, and work is neither applied nor extracted unless an additional external non-conservative force is applied. In the transformed coordinates, the picture is slightly different, as there is only a single heat bath. To find the heat and work rates, we derive the steady state distribution of this system. This can be done explicitly by directly solving the Klein-Kramers equation for the curl-force OU process \cite{risken1996fokker} using computer algebra \cite{gowda2021high}, or by simply transforming the known solution of the underdamped Brownian gyrator \cite{mancois2018two, bae2021inertial} to the curl-force coordinates. The solution is a Gaussian $p(\mathbf{z}) \sim \exp(-\frac{1}{2}\mathbf{z}^T C^{-1} \mathbf{z})$, where $\mathbf{z} = (x, y, \dot{x}, \dot{y})^T$ and $C$ is the steady-state covariance matrix. The energetics of the system are computed directly from $C$. From here on, we fix the temperature to $T=\sqrt{T_X T_Y}$ (and correspondingly $\alpha\gamma=1$) to simplify the expressions. Other choices of $T$ are possible, and the results are dependent on this choice. The power applied on the system by the non-conservative force is
\begin{equation}\label{eq:work}
    \langle \dot{W} \rangle = \langle \mathbf{F} \cdot \dot{\mathbf{x}} \rangle = 
\frac{k_B T\xi u^2(\alpha-\gamma)^2 }{2  \xi^{2}  k + 2  u^{2}  m }.
\end{equation}
We note that $\langle \dot{W}\rangle \sim |\nabla \times \mathbf{F}|^2$, such that the curl of the force field is a measure of how far the system is from thermal equilibrium, in the sense that a larger curl implies that more work is applied on the system. 
We note that Eq.~\ref{eq:work} also implies $\langle \dot{W}\rangle\geq0$: positive work is always applied on the system regardless of the direction of gyration (determined by whether $\alpha > \gamma$ or $\alpha < \gamma$), which is a manifestation the second law: work cannot be extracted from a single heat bath. This emphasizes the difference of this description from the two baths description, where heat flow is from the hot bath to the cold one \cite{bae2021inertial}. In that case, the direction of the flow can be reversed by interchanging the bath temperatures.

Given the solution for $C$, the average heat dissipation rate is also computed directly as $\langle \dot Q \rangle= -\frac{2\xi}{m}\left(\frac{m}{2}(\langle \dot{x}^2\rangle + \langle \dot{y}^2\rangle) - k_B T\right)$ \cite{sekimoto_stochastic_2010}. Since  the system is in a steady state, the work done by the non-conservative force is fully dissipated into the bath as heat, implying also that \begin{equation}\label{eq:heat}
\langle \dot{Q} \rangle=-\langle \dot{W}\rangle \leq 0.
\end{equation}
Finally, this allows the computation of the entropy production rate
\begin{equation}
    \langle \dot{S}\rangle=\langle \dot{S}\rangle_\text{system}+\langle \dot{S}\rangle_\text{bath}=\langle \dot{S} \rangle_\text{bath}=-\frac{\langle \dot{Q}\rangle}{T} = \frac{k_B \xi u^2(\alpha-\gamma)^2 }{2  \xi^{2}  k + 2  u^{2}  m } \geq 0.
\end{equation}
In terms of the two-bath system parameters, substituting $\alpha=\sqrt{T_Y/T_X}$ and $\gamma=\sqrt{T_X/T_Y}$, we get
\begin{equation}\label{eq:gyrator_EPR}
    \langle \dot{S}\rangle = \frac{k_B \xi u^2(T_Y-T_X)^2 }{2T_X T_Y(  \xi^{2}  k +  u^{2}  m)},
\end{equation}
which is identical to the EPR obtained for the underdamped gyrator \cite{bae2021inertial}, as expected \cite{godreche2019characterising}: unlike quantities like heat or work, which measure the flow of energy between the different components of the full system (bath, non-conservative force, the system itself), the total entropy production rate measures the irreversibility of the system as a whole, and does not depend on the representation we use to describe it. This reflects the fact that the NESS of the gyrator and the curl-force system are related by a linear coordinate transformation \cite{godreche2019characterising}.

Finally, as a sanity check we also computed the EPR using different choices for the temperature, taking $T$ as the weighted arithmetic or geometric mean of $T_X$ and $T_Y$, $T=wT_X+(1-w)T_Y$ or $T=T_X^wT_Y^{1-w}$ respectively (for $0 < w < 1$). In both cases, the same EPR as in Eq.~\ref{eq:gyrator_EPR} is also obtained. Rates of work and heat, on the other hand, do depend on this choice, because they explicitly depend on the bath's temperature.
A choice that can be of particular interest is to require that $\langle \dot W\rangle = \langle \dot Q _X \rangle$ and $\langle \dot{Q}\rangle = \langle \dot{Q}_Y\rangle$ (for $T_X>T_Y$), such that the flows of energy into or out of the system is the same in the two descriptions. An exact solution can be obtained by comparing the expressions for the heat flow $\langle \dot{Q} \rangle = -T\langle\dot S\rangle$ with that obtained for the gyrator \cite{bae2021inertial}, which yields the relation
\begin{align}
    T = \begin{cases}
                T_{X, Y} & T_X=T_Y \\
            \frac{T_X T_Y}{|T_X - T_Y|} & T_X \neq T_Y
    \end{cases}.
\end{align}
We do not know whether this is a general result, or if it is system dependent. Notably, if $T_X \gg T_Y$ (or $T_Y \gg T_X$), then the single bath's temperature for which the heat flows are equal to those in the multi-bath system is approximated by the temperature of the colder bath.

\subsection{Beyond the quadratic potential: gyrator on a ring}
To illustrate the generality of this framework, we consider the gyration of a particle confined to a ring by the potential $\tilde{U}(R) = \frac{1}{2}k(R-R_0)^2$, where $R=\sqrt{X^2+Y^2}$ is the radial distance from the center of the ring and $R_0$ is its radius \cite{abdoli2025quadrupolar}. The system is coupled to two orthogonal sources of thermal noise with temperatures $T_{X}$ and $T_{Y}$, as in the case of the harmonically confined gyrator. The flux field in this system has been shown to be of a quadrupolar structure, with four alternating regions of clockwise and counter-clockwise rotating currents \cite{abdoli2025quadrupolar}. Transforming the system to the curl-force representation, we obtain the force field
\begin{equation}
   \mathbf{F}(\mathbf{x}) \equiv -\mathcal{A}\cdot\nabla_{\mathbf{x}}U = -k\left( 1 - \frac{R_0}{\sqrt{\frac{x^2}{\alpha} + \frac{y^2}{\gamma}}} \right) \mathbf{x} = -k\left( 1 - \frac{\sqrt{T}R_0}{\sqrt{T_X{x^2} + T_Yy^2}} \right) \mathbf{x}
\end{equation}
where $\alpha = T/T_X$ and $\gamma = T/T_Y$ with $T$ the temperature of the transformed, single-reservoir system. This force field is radial and is confining to the ring, which is the ellipse ${\frac{x^2}{\alpha} + \frac{y^2}{\gamma}} = R_0^2$ in the dual system. Fig.~\ref{fig:curl_ring_gyrator}a shows the force field and the value of its curl,
\begin{equation}
    (\nabla \times \mathbf{F})\cdot\hat{z} = \frac{R_0 k \left(\alpha - \gamma\right) x y}{\alpha\gamma\left( \frac{x^2}{\alpha} + \frac{y^2}{\gamma} \right)^{3/2}} = \frac{R_0 k \sqrt{T} (T_Y - T_X)xy}{(T_X x^2 + T_Y y^2)^{3/2}}.
\end{equation}
Unlike the case of the harmonic gyrator, the curl of the force field is not constant, and has a quadrupolar structure instead. For $T_X>T_Y$ ($\alpha<\gamma$), it is negative (positive) in the 1st and 3rd (2nd, 4th) quadrants, corresponding to clockwise (counter-clockwise) circulation. Again, the direction of circulation predicted by the sign of $\nabla \times \mathbf{F}$ coincides with those calculated for the gyrator on a ring \cite{abdoli2025quadrupolar}. We numerically verify that the same direction is also observed in the curl-force system by performing Langevin simulations \cite{rackauckas2017differentialequations, rackauckas2017adaptive, rossler2010runge} of the dynamics, as shown in Fig.~\ref{fig:curl_ring_gyrator}b. Unlike the case of the harmonic Brownian gyrator, in this case the steady-state solution has no exact analytical solution, and requires approximations or numerical solutions. Hence, in this case, the easy exact computation of $\nabla \times \mathbf{F}$ provides a simple method for obtaining some qualitative properties of the system, which are otherwise not directly accessible.

\begin{figure}
    \centering
    \includegraphics[width=\linewidth]{velocity_field.pdf}
    \caption{(a) The force field $\mathbf{F}=m\ddot{\mathbf{x}}$ of the single-bath system dual to the ring gyrator (arrows), and the value of its curl $\nabla\times\mathbf{F}$ (color gradient), for $k = 20, R_0 = 1, \alpha = \sqrt{3/2}$ and $\gamma = \sqrt{2/3}$. Unlike the case of the harmonically confined Brownian gyrator, here $\nabla \times \mathbf{F}$ is not constant, and has a quadrupolar structure instead.
    (b) Simulation results for the probability density and currents averaged over 25,000 Langevin trajectories of the same system coupled to a bath with temperature $k_BT=\sqrt{3/2}$ and drag coefficient $\xi=1$ ($k_B = 1$). This temperature corresponds to $T=\sqrt{T_Y T_X}$ with $T_X= 1$ and $T_Y=3/2$, matching the choice of $\alpha, \gamma$ in the left panel.
    The steady state probability density function is shown in the background (white - low density, black - high density). The arrows indicate the empirical spatial probability current $\mathbf{J}_\mathbf{x}=\langle \mathbf{v} | \mathbf{x}\rangle f$, where $\langle \mathbf{v} | \mathbf{x}\rangle$ is the average velocity conditioned on position.
    A positive value of $\nabla\times\mathbf{F}$ in (a) corresponds to counter-clockwise rotation in (b) and a negative value corresponds to clockwise rotation. These directions are in agreement with those found for the gyrator on a ring \cite{abdoli2025quadrupolar} for the case $T_Y>T_X$ (corresponding $\alpha>\gamma$ as evaluated here).
    The ring $X^2 + Y^2 = R_0^2$, which is transformed into an ellipse ${\frac{x^2}{\alpha} + \frac{y^2}{\gamma}} =R_0^2$ in the dual system, is shown in a solid line in both figures.}
    \label{fig:curl_ring_gyrator}
\end{figure}
\section{Discussion and outlook}\label{sec:discussion}
In this work, we studied systems that are driven out of equilibrium by a force which is non-conservative but can nevertheless be described by a Hamiltonian with a non-standard form. First, by assuming that the non-conservative forces do not modify the system's interaction with the heat bath, we have shown that such systems are generally not expected to relax to equilibrium despite their Hamiltonian structure.
We then presented a mapping between two distinct mechanisms that drive a system out of equilibrium: coupling different DoF to heat baths at different temperatures, and the presence of a specific type of Hamiltonian curl-forces. This allows for a mechanical interpretation of the origin of rotational motion often observed in systems coupled to multiple heat baths, such as the Brownian gyrator, as the result of deterministic non-conservative forces. The direction of gyration is qualitatively described by the curl of this force field, $\nabla\times\mathbf{F}$, which is obtained without solving the stochastic dynamics of the system. Such interpretation has the benefit of being intuitive and easier to grasp in comparison to alternative ways of qualitatively explaining such phenomena, such as through an analysis of the angular dependence of the diffusion tensor \cite{abdoli2025quadrupolar}. It is clear that $\nabla\times\mathbf{F}$ cannot fully describe the currents in the thermal system as it lacks information regarding the steady state distribution itself. However, we note that other quantities suggested for the characterization of the non-trivial dynamics arising in such systems, such as the curl of the probability current \cite{dotsenko2013two}, are not necessarily more informative, as they tend to be complicated and do not carry a clear physical significance. Hence, in some cases $\nabla \times \mathbf{F}$, which is easily computed for arbitrarily complex systems, may be sufficient to describe some properties of the system. We further note that if $\nabla \times \mathbf{F} = 0$ in the curl-force representation, then the system is in thermal equilibrium with the bath. In the multiple-bath description, this is reflected in each degree of freedom independently equilibrating with its own bath.

The exactly solvable case of the OU process allowed us to explicitly compute the energetics of the system with curl-forces. In the steady state, we found that positive work is always applied on the system, regardless of the direction of gyration. This work is fully dissipated into the single heat bath as heat, and is associated with a positive EPR. Notably, while the mechanism driving the system out of equilibrium in each system is different (heat flowing between two reservoirs, compared to work being dissipated in the curl-force system), the EPR\textcolor{blue}{---an information theoretic quantity---}of the two systems is equal. Nevertheless, the two systems are distinguishable from one another, because experimentally measurable quantities such as the dissipated heat are not the same in the two descriptions.
We note that in contrast to gyrators, where heat flows from the bath with the higher temperature to that with the lower one, in the curl-force system work is always applied on the system and dissipated into the single bath, a consequence of the second law of thermodynamics.

% A powerful implication of the presented transformation is the possibility of making non-trivial qualitative predictions on systems for which an exact solution is unknown, as discussed in the example of the gyrator on a ring. In principle, the same analysis of $\nabla\times \mathbf{F}$ can be applied to other, arbitrarily complicated force fields, $\mathbf F$, to gain insights on the currents expected in a NESS (if such exists).

More generally, it will be interesting to further investigate the role of $\nabla \times \mathbf{F}$ on the nonequilibrium dynamics and energetics of a driven system. We have shown that not only it can be used to predict the direction of circulating currents in a NESS in both cases we studied, but also that in the exactly solvable linear case, where $\nabla\times\mathbf{F}$ is a constant, quantities such as power, heat dissipation rate and EPR are a function of its magnitude. Whether and how this result generalizes to other systems remains an open question.

One especially promising application of our method may be in experimental realizations of systems with multiple heat baths. Some experiments have been performed in the past, for example using system analog to Brownian motion consisting of an electric circuit with resistors held at different temperatures \cite{ciliberto2013heat, chiang2017electrical}. There have also been experiments conducted directly on Brownian particles, where in this case the anisotropic noise is artificially implemented by an athermal random force \cite{argun2017experimental, berut2014energy, lin2022stochastic, berut2016stationary, martins2026role}. However, none of these experiments featured diffusion with ``real'' thermal anisotropic noise, because it is not clear what it means for a single Brownian particle to be coupled to two heat sources simultaneously. On the contrary, curl-forces may be readily generated in the lab with existing setups, for example using optical tweezers \cite{roichman2008influence,roichman2008optical, sun2009brownian, wu2009direct, liberal2013near, sukhov2017non}. This could pave the way for an easier approach to effectively simulating experiments with multiple temperatures by performing experiments on the curl-force description instead, before transforming results back to the multiple-temperatures description.

This work was limited to a small subset of curl-forces that are described by a Hamiltonian with anisotropic---yet quadratic and positive---kinetic energy \cite{berry2015hamiltonian}. Generalizing to other types of non-conservative forces or Hamiltonians \cite{yip2024general}, as well as non-mechanical thermodynamic forces, could possibly lead to mappings of other types of nonequilibrium systems beyond the gyrator and its immediate generalizations. In essence, applying a similar canonical transformation to arbitrary nonequilibrium stochastic systems could lead to a more general form of ``nonequilibrium Hamiltonians'', which could eventually help form a unified theory for nonequilibrium systems \cite{ding2022unified}. One promising direction in this respect is performing a more rigorous mathematical study of curl-force dynamics and the identification of new conserved quantities, for example as was recently done by Yavari and Goriely \cite{yavari2025darboux}. While the conserved quantity they identified (specifically, a Hamiltonian) is only defined implicitly in terms of an integral equation so it cannot be applied directly to any system with curl-forces, perhaps explicit solutions could be derived for specific systems. In combination with our method, this could be a step towards finding closed form NESS solution for certain systems with multiple heat baths beyond the exactly solvable OU process. Finally, it is interesting to note that the Brownian gyrator can also be mapped to an information engine \cite{leighton2024information}. It is worthwhile considering how the curl-force description of the system relates to the feedback-control one, and more generally how these two methods of driving a system out of equilibrium relate to one another.

Other extensions of this work are also possible. It is interesting to consider whether this framework can be extended to the overdamped regime. The Brownian gyrator was originally formulated in the this limit \cite{filliger2007brownian}, but within our framework we are unable to give a mechanical interpretation of the rotational motion in this limit, as our description relies on the definitions of momentum and kinetic energy. Another possible application is studying the curl-force representation of gyrators with potentials other than the harmonic or ring potentials studied here, such as those considered in \cite{chang2021autonomous}, which could possibly provide more general insights. Finally, it is intriguing to construct a kinetic theory for a closed system consisting of particles with a non-isotropic kinetic energy term, namely when the bath itself has a curl-force Hamiltonian. Would such a system relax to the corresponding equilibrium, or to a nonequilibrium steady state? 

 \paragraph{Acknowledgments}
\textcolor{red}{Add Alja\v{z}'s funding and Omer's (if any should be listed)}
OR acknowledges financial support from the ISF, grant no. 232/23,  by the Minerva foundation with funding from the Federal German Ministry for Education and Research and by the Lower Saxony - Israel collaboration grant. We would like to thank Roi Holtzman and Shahaf Aharony Shapira for helpful discussions.
 
% \bibliographystyle{RS}
% \bibliography{main}
\printbibliography

\end{document}